% XVI-P p19 note: the projected one-loop kappa^2 ch-identity defect
\section{P.5 identity step (p19): the projected one-loop $\kappa^2$ ch-identity defect}
\scopebox{Scope.}{p19 keeps the p14r contract and the p15--p18 calculation layers byte-identical
(\texttt{release\_p19.py} recomputes their sha256). It assembles the ch projection (P.3, frozen) of the
$\hbar\kappa^2$ part of $(S_\Psi,\Gamma_\Psi)$ from every evaluated input, exactly in $d$. It then solves for a
symmetry-restoring counterterm in the declared domain. Only the ch projection is covered (one $c$, one $h$; one loop;
$\kappa^2$; off shell). bulk\_symmetry\_restored is false. O.1 stays NOT\_CLOSED; the insertion stays frozen; XV-B
and XIV-J are unopened; Route B is not discharged; no $c\,\delta$ is added to $u^R_\mu$.}

\subsection{Assembly}
The ledger \texttt{P13\_ch\_dependency\_rows.json} has three surviving rows:
\begin{itemize}
\item $R_{hh}$, $(S_0,\Gamma_1^{[2]})$ $s_2K{:}K.c$, inputs D01+D02+D03+D11+D12, with
  $X_{hh}=\int(s_0h)\,G_{hh}\,h$;
\item $R_{Kc}$, $(S_0,\Gamma_1^{[2]})$ $s_1h{:}h^2$, input D13, with $X_{Kc}=\int G_{Kc}\,c\,E\,h$, where $E$ is
  the $h^2$-monomial operator of $S_0$;
\item $R_{1pt}$, $(S_1,\Gamma_1^{[1]})$, inputs D14a--c.
\end{itemize}
Every input is evaluated; none is Missing, and a Missing input is refused rather than set to zero (control). The
defect is computed for seven external assignments, including ``c hat, h tilde''. For each it is given as the exact
$d$-dependent value times MB, its pole, its coefficient of $L$ and its MS-bar local part.

\paragraph{Result: the defect is zero, exactly in $d$, in every part and every assignment.}
\textbf{Why it is zero.} Each surviving row vanishes separately in the projected slot. No counterterm cancellation
is involved, and no cancellation between rows.
\begin{itemize}
\item[(i)] hh row: full-row transverse decoupling. D01 and D03 are each nonzero in the slot and cancel each other
  (the p17/p18 transverse sum). D11 is transverse by itself, and D02 = D12 = 0.
\item[(ii)] K.c row: D13 reaches the ch slot only through $T_1$ and $E_1$ ($T_2$, $T_3$ and $E_2$ are longitudinal;
  $E$ is transverse), and $A_{T_1}=A_{E_1}=0$ exactly in $d$.
\item[(iii)] One-point row: the tadpoles D14a--c vanish under the declared prescription.
\end{itemize}
\textbf{This rests on:}
\begin{itemize}
\item M2 / the D11 matching, verified on the \emph{4D components only} (the D11 matching for the evanescent
  components is not applicable and not verified);
\item the evanescent decoupling proofs (P16 E2; P17/P18 hh evanescent and full-row decoupling);
\item the frozen P.3 projection and the P13 ledger;
\item the declared prescription.
\end{itemize}
\textbf{Evanescent decoupling in the assembly.} The evanescent parts of both rows were assembled explicitly and are
zero for all seven assignments. The evanescent poles themselves are nonzero; they are regulated and not subtracted.

\subsection{Counterterm domain and solve}
Route G applies $\sigma_0=(S_0,\cdot)$ to the 11 p14r domain densities: h.h\#0--4, c.cb\#0--1, c.K\#0--2,
$\phi.\phi$. Route C computes the covariant ch images exactly in $d$. Both give rank $P_{ch}\sigma_0=3$.
\begin{itemize}
\item Over $\mathbb Q(i)$ (these counts are the real dimensions under a compatible real structure; wording corrected
  in p20): rank $\sigma_0|_D=5$, $\dim D\cap\ker\sigma_0=6=\dim D\cap{\rm Im}\,
  \sigma_0$, cohomology 0, and the quotient has dimension 5 (complement h.h\#0--2, c.cb\#0--1).
\item The h.h + c.K $\Gamma_1$ data, without the ghost self-energy and in route-G (c,cb) labels, have quotient
  coordinates pole $(0,0,0,-3/8,1/4)$, coefficient of $L$
  $(0,0,0,3/8,-1/4)$ and local part $(0,0,0,-2/3,5/12)$. The h.h pole is transverse, i.e.\ $\sigma_0$-exact. These coordinates equal $-[\Sigma]$: $\Sigma$ itself has
  $(+3/8,-1/4)$, and the data including $\Sigma$ have quotient $0$ [clarified in p20].
\item The linear solve $\Delta+P_{ch}\sigma_0(x)=0$ gives $x=0$ (outcome ZERO); no residual and no obstruction.
\item The MS-bar counterterm has zero ch image. The h.h + c.K part alone has a nonzero full $\sigma_0$ image in the
  c.b sector; adding the c.cb ghost self-energy counterterm in the correct orientation closes it [corrected in p20,
  see the p20 erratum note].
\end{itemize}

\paragraph{Flags and their definitions.}
\begin{itemize}
\item \texttt{projected\_defect\_computed}: $\Delta$ is assembled from all evaluated inputs, exactly in $d$, in
  all parts and all seven assignments.
\item \texttt{projected\_identity\_verified}: $\Delta=0$ for the \emph{regulated} $\Gamma_1$.
\item \texttt{projected\_restoration\_verified} [REDEFINED in p20]: full $\sigma_0$ closure, in every ghost $+1$
  sector (h.c, c.b, c.cbs), of the \emph{complete} one-loop $\kappa^2$ MS-bar counterterm $S_{ct}$ (h.h + c.K + c.cb,
  ghost self-energy held fixed, derived relative factor $+1$, orientation converted). The p19 definition
  ($P_{ch}\sigma_0(S_{ct})=0$, $x=0$) follows from the identity and locality; it is kept only as a labelled
  route-consistency control.
\end{itemize}
\emph{Scope of the restoration flag:} the one-loop $\kappa^2$ counterterm only; not bulk symmetry, not O.1. The
release gate re-checks each definition from the certificate content. bulk\_symmetry\_restored is false regardless.

\paragraph{K1 [corrected in p20; the p19 text is retracted].} p19 reported that full closure of the MS-bar
counterterm fails with the derived ghost-equation factor $+1$ and needs $-1$. That was a Grassmann-orientation
bookkeeping error: the route-G densities are built in the order (c,cb), whereas the P15 $\Sigma$ coordinates are
(cb,c)-oriented, and $V[\mathrm{c.cb}\#k]=-V[(\mathrm{cb,c})\#k]$. With the measured ratio $-1$ applied, the derived
factor $+1$ gives full closure (p20 erratum note). The relations below are stated in the (cb,c) orientation:
\begin{itemize}
\item The ghost-equation operator $2iD$, verified by P15 as $\Sigma=2iD[G_{Kc}]$, gives
  $2iD\,V[\mathrm{c.K\#1}]=-V[(\mathrm{cb,c})\#0]$ and $2iD\,V[\mathrm{c.K\#2}]=-V[(\mathrm{cb,c})\#1]$.
\item Applied to the c.K pole, this reproduces the P15 $\Sigma$ pole $(-3/8,1/4)$ in (cb,c) coordinates, i.e.\
  $(+3/8,-1/4)$ in the route-G (c,cb) coordinates.
\end{itemize}

\paragraph{J1: finding about the p14r contract (contract not modified).} The p14r counterterm-domain certificate
counts Im $\sigma_0$ 3, quotient 8, $D\cap\ker$ 5 and rank $P_{ch}\sigma_0$ 4. These are artefacts of a real/imaginary
split over $\mathbb Q$: the densities carry an explicit $i$ for odd derivative counts, so an $i$-multiple counts as
independent. Over $\mathbb Q(i)$ the counts are 6, 5, 6 and 3 (the real dimensions under a compatible real
structure), and Route C agrees. With the split, p19 reproduces the
p14r counts (control).

\paragraph{Controls and mutations.} All controls are true. Eight mutations each change the outcome:
\begin{itemize}
\item drop D03; drop D01; flip the D03 phase; double D03; trace operator instead of $s_0h$; de Donder $q^2P$ instead
  of $E$; drop the D03 evanescent part (naive 4D). Each gives OBSTRUCTION (nonlocal residual).
\item D13 with $G_t$ dropped gives NONZERO but CANCELLABLE (rank 3).
\end{itemize}
Recorded blind spots of the projected identity: dropping D11 (transverse by itself), dropping D13
($A_{T_1}=A_{E_1}=0$) and flipping the K.c row sign (each row vanishes separately). In p20, the full-closure test
detects the last two, and the D11 normalisation pin covers the first. Removing D11 from the inventory makes the inventory incomplete, and
projected\_defect\_computed is then not earned.

% generated by code/p5_calc/release_p19.py -- do not edit
\begin{center}\scriptsize\begin{tabular}{p{3.2cm}p{2.6cm}p{2.6cm}p{2.6cm}}\hline
assignment $(c,h)$ & defect pole & coeff. of L & local finite \\\hline
\texttt{Q4S (declared, R2)} & \texttt{0} & \texttt{0} & \texttt{0} \\
\texttt{QdS (hat)} & \texttt{0} & \texttt{0} & \texttt{0} \\
\texttt{c tilde, h hat} & \texttt{0} & \texttt{0} & \texttt{0} \\
\texttt{c tilde, h mixed} & \texttt{0} & \texttt{0} & \texttt{0} \\
\texttt{c hat, h mixed} & \texttt{0} & \texttt{0} & \texttt{0} \\
\texttt{all tilde} & \texttt{0} & \texttt{0} & \texttt{0} \\
\texttt{c hat, h tilde} & \texttt{0} & \texttt{0} & \texttt{0} \\
\hline\end{tabular}\end{center}
\noindent Route G over \texttt{Q(i)}: rank $\sigma_0|_D$ = 5, dim $D\cap\ker\sigma_0$ = 6, dim $D\cap\mathrm{Im}\,\sigma_0$ = 6, cohomology 0, quotient 5 (complement \texttt{h.h\#0, h.h\#1, h.h\#2, c.cb\#0, c.cb\#1}); rank $P_{ch}\sigma_0$ = 3 (Route C: 3).
Over Q with the re/im split (p14r convention): dim $D\cap\mathrm{Im}$ = 3, quotient 8, dim $D\cap\ker$ = 5.

\noindent Quotient coordinates of the $\Gamma_1$ local data: pole \texttt{['0', '0', '0', '-3/8', '1/4']}, coeff. of L \texttt{['0', '0', '0', '3/8', '-1/4']}, local \texttt{['0', '0', '0', '-2/3', '5/12']}.
\begin{center}\scriptsize\begin{tabular}{p{4.6cm}p{6.4cm}p{1.2cm}}\hline
mutation & outcome & changes \\\hline
\texttt{drop\_graph\_D03} & \texttt{OBSTRUCTION (nonlocal residual: the L-coeffi} & yes \\
\texttt{drop\_graph\_D01} & \texttt{OBSTRUCTION (nonlocal residual: the L-coeffi} & yes \\
\texttt{flip\_phase\_D03} & \texttt{OBSTRUCTION (nonlocal residual: the L-coeffi} & yes \\
\texttt{symmetry\_factor\_D03\_x2} & \texttt{OBSTRUCTION (nonlocal residual: the L-coeffi} & yes \\
\texttt{wrong\_projection\_trace\_operator\_hh} & \texttt{OBSTRUCTION (nonlocal residual: the L-coeffi} & yes \\
\texttt{wrong\_projection\_gauge\_fixed\_operator\_Kc} & \texttt{OBSTRUCTION (nonlocal residual: the L-coeffi} & yes \\
\texttt{naive\_4D\_D03\_evanescent\_dropped} & \texttt{OBSTRUCTION (nonlocal residual: the L-coeffi} & yes \\
\texttt{naive\_4D\_D13\_Gt\_dropped} & \texttt{NONZERO, CANCELLABLE by a counterterm in the} & yes \\
\hline\end{tabular}\end{center}
\noindent Counterterm solve: \texttt{ZERO}; \texttt{x = 0 (no counterterm needed beyond the MS-bar poles); the identity-preserving finite renormalisations are ker(P\_ch sigma\_0) cap D (dim 8), of which 2 modulo D cap Im sigma\_0}.

\noindent\textbf{Why zero.} \texttt{each surviving ch row vanishes separately in the projected slot; no counterterm cancellation and no cancellation between rows is involved: (i) hh row (D01+D02+D03+D11+D12): full-row transverse decoupling -- D01 and D03 are individually nonzero and cancel against each other in the slot (p17/p18 transverse sum), D11 is transverse by itself, D02 = D12 = 0; (ii) K.c row (D13): D13 reaches the ch slot only through T1 and E1 (T2, T3, E2 are longitudinal, E transverse) and A\_T1 = A\_E1 = 0 exactly in d; (iii) one-point row: the tadpoles D14a-c vanish under the declared prescription.  The defect is therefore zero without any counterterm (solve x = 0).}

\noindent\textbf{Flag definitions.} \textbf{\texttt{projected\_defect\_computed}}: \texttt{Delta = P\_ch[(S\_Psi, Gamma\_Psi)]\_\{hbar kappa\^{}2\} assembled from every evaluated input of the three surviving rows (none Missing), exact in d, pole + nonlocal + MS-bar local parts, all seven external assignments}. \textbf{\texttt{projected\_identity\_verified}}: \texttt{Delta = 0 for the REGULATED one-loop functional Gamma1: pole, nonlocal and local parts vanish for all seven assignments, exactly in d, with all controls true and every mutation changing the outcome}. \textbf{\texttt{projected\_restoration\_verified}}: \texttt{p20 definition (replaces p19's, which follows from the identity + locality): FULL sigma\_0 closure of the COMPLETE one-loop kappa\^{}2 MS-bar counterterm S\_ct = -(pole) in the declared domain, h.h + c.K + c.cb (ghost self-energy Sigma held fixed at its P15 value, orientation converted, derived relative factor +1), in every ghost +1 sector (h.c, c.b, c.cbs), together with projected\_identity\_verified.  Scope: the one-loop kappa\^{}2 counterterm only; NOT bulk symmetry, NOT O.1}. \textbf{\texttt{bulk\_symmetry\_restored}}: \texttt{hard stop: False regardless (not claimed; O.1 NOT\_CLOSED)}.

\noindent\textbf{K1 (corrected in p20; erratum to p19).} derived relative factor \texttt{1}; orientation ratio (route-G (c,cb) over P15 (cb,c)) \texttt{['-1', '-1']}; Sigma pole in (c,cb) coordinates \texttt{['3/8', '-1/4']}; full sigma\_0 image of the complete MS-bar counterterm: \texttt{\{'closed': True, 'entries': 0, 'sectors': []\}}. \texttt{ERRATUM to p19: the p19 statement 'closure fails with the derived factor, needs -1, unresolved' was a Grassmann-orientation bookkeeping error (P15 (cb,c) coordinates placed in route-G (c,cb) slots).  Corrected: derived +1, full closure holds (finding K1 of the p19 independent rebuild).}

\noindent\textbf{J1 (finding about the p14r contract).} \texttt{FINDING ABOUT THE p14r CONTRACT (contract NOT modified): the p14r counterterm-domain counts (Im sigma\_0 3, quotient 8, D cap ker 5, P\_ch sigma\_0 rank 4) are artefacts of a re/im split over Q (the densities carry explicit i for odd derivative counts, so i-multiples count as independent).  The counts over Q(i) -- sigma\_0 rank 5, D cap ker 6, D cap Im 6, cohomology 0, quotient 5, P\_ch sigma\_0 rank 3 (= Route C) -- equal the real dimensions under a compatible real structure.  The p14r certificate stays as released; both are reported}

